\section{Objective}
\begin{itemize}[roomy]
	\item To understand DeMorgan's Theorems and their role in Boolean algebra and digital logic design.
	\item To verify DeMorgan's Theorems experimentally for NOT, AND and OR gates using the 7404, 7408 and 7432 logic gate ICs.
	\item To confirm from the recorded truth tables that $\overline{A+B}=\nA\cdot\nB$ and $\overline{A\cdot B}=\nA+\nB$.
\end{itemize}

\section{Introduction}
Digital circuits are built from a handful of basic gates, and a lot of design work comes down to rewriting an expression in a simpler or more convenient form. \textbf{DeMorgan's Theorems} are two of the most useful rules for doing that: they show how an AND operation can be turned into an OR operation (and the other way round) just by inverting the inputs and the output. This is also why a circuit can be built using only NAND or NOR gates. In this experiment we build an inverter circuit and two two-LED gate circuits on a breadboard using TTL logic gate ICs, drive them with push buttons, and compare the LED outputs with what the two theorems predict.

\section{Theory}

\subsection{DeMorgan's First Theorem}
\begin{quote}\itshape
	The complement of a logical product (AND) is equal to the logical sum (OR) of the complements of the individual variables.
\end{quote}
\[
\overline{A\cdot B}=\nA+\nB
\]
In gate terms, a NAND gate is equivalent to an OR gate with inverted inputs (a negative-OR or ``bubbled OR'' gate).

\subsection{DeMorgan's Second Theorem}
\begin{quote}\itshape
	The complement of a logical sum (OR) is equal to the logical product (AND) of the complements of the individual variables.
\end{quote}
\[
\overline{A+B}=\nA\cdot\nB
\]
In gate terms, a NOR gate is equivalent to an AND gate with inverted inputs (a negative-AND or ``bubbled AND'' gate).

\subsection{Extension to More Variables}
Both theorems extend to any number of variables, for example
\[
\overline{A\cdot B\cdot C}=\nA+\nB+\overline{C},\qquad \overline{A+B+C}=\nA\cdot\nB\cdot\overline{C}.
\]

\subsection{Logic Gate ICs Used}
\begin{itemize}
	\item \textbf{7404} -- hex inverter, six independent NOT gates in one package.
	\item \textbf{7408} -- quad 2-input AND gate, four independent AND gates in one package.
	\item \textbf{7432} -- quad 2-input OR gate, four independent OR gates in one package.
\end{itemize}
All three are 14-pin TTL ICs powered from a 5\,V supply, with $V_{cc}$ on pin 14 and GND on pin 7. An LED with a series current-limiting resistor connected to a gate output shows the logic level: a HIGH output (logic 1) turns the LED on and a LOW output (logic 0) leaves it off.

\section{Apparatus and Components}
\begin{table}[H]
	\centering
	\caption{List of Equipment and Components Used}\label{tab:apparatus10}
	\small
	\begin{tabular}{@{}llc@{}}
		\toprule
		Item Description & Specification / Rating & Quantity \\
		\midrule
		Breadboard & Standard Tie-Point & 1 \\
		Hex Inverter IC (NOT gates) & 7404 & 1 \\
		Quad 2-Input AND Gate IC & 7408 & 1 \\
		Quad 2-Input OR Gate IC & 7432 & 1 \\
		LEDs & Different colours for easy identification & 2 \\
		Current-Limiting Resistors & 220\,$\Omega$ or 330\,$\Omega$ & 2 \\
		Push Buttons (SW1, SW2) & Momentary & 2 \\
		DC Power Supply & 5\,V & 1 \\
		Connecting Wires & Jumper Wires & As required \\
		\bottomrule
	\end{tabular}
\end{table}

\section{Circuit Diagrams}
\begin{figure}[H]
	\centering
	\includegraphics[width=0.58\textwidth]{circuits/ic7404.png}
	\caption{NOT Gate IC 7404 (hex inverter) pin diagram, used for the inverter circuit and to generate $\overline{A}$ and $\overline{B}$.}
	\label{fig:ic7404}
\end{figure}

\begin{figure}[H]
	\centering
	\includegraphics[width=0.58\textwidth]{circuits/ic7408.png}
	\caption{AND Gate IC 7408 (quad 2-input AND) pin diagram.}
	\label{fig:ic7408}
\end{figure}

\begin{figure}[H]
	\centering
	\includegraphics[width=0.58\textwidth]{circuits/ic7432.png}
	\caption{OR Gate IC 7432 (quad 2-input OR) pin diagram.}
	\label{fig:ic7432}
\end{figure}

\section{Experimental Procedure}
\textbf{Part A: Inverter (NOT gate)}
\begin{enumerate}
	\item Place the 7404 on the breadboard and connect $V_{cc}$ (pin 14) to $+5$\,V and GND (pin 7) to ground.
	\item Connect push button SW1 to the input of one inverter, and connect an LED with a series resistor from that inverter's output to ground (Figure~\ref{fig:ic7404} shows the 7404 pinout).
	\item Press and hold SW1 (input LOW) and note the LED state; release it (input HIGH) and note the LED state. Record both in the truth table.
\end{enumerate}
\textbf{Part B: DeMorgan's Theorem with AND gates}
\begin{enumerate}[resume]
	\item Place the 7408 on the breadboard and power it. Connect SW1 and SW2 to the two inputs of one AND gate and its output to LED~1 through a resistor.
	\item Pass SW1 and SW2 through two inverters of the 7404 and feed the inverted signals $\nA$ and $\nB$ to a second AND gate, whose output drives LED~2 (7408 pinout in Figure~\ref{fig:ic7408}).
	\item Operate SW1 and SW2 through all four input combinations and record the inverted inputs and the state of LED~1 and LED~2 each time.
\end{enumerate}
\textbf{Part C: DeMorgan's Theorem with OR gates}
\begin{enumerate}[resume]
	\item Replace the two AND gates with OR gates from the 7432 and keep the rest of the wiring the same (7432 pinout in Figure~\ref{fig:ic7432}), so that LED~1 shows $A+B$ and LED~2 shows $\nA+\nB$.
	\item Repeat all four input combinations and record the results in a new truth table.
\end{enumerate}

\section{Theoretical Analysis}
The expected behaviour follows directly from the two theorems. Table~\ref{tab:theory10} lists every term for all four input combinations. The columns $\overline{A+B}$ and $\nA\cdot\nB$ are identical, and so are $\overline{A\cdot B}$ and $\nA+\nB$, which proves the theorems.

\begin{table}[H]
	\centering
	\caption{Expected Truth Table for DeMorgan's Theorems}\label{tab:theory10}
	\small
	\setlength{\tabcolsep}{5pt}
	\begin{tabular}{@{}cccc|ccc|ccc@{}}
		\toprule
		$A$ & $B$ & $\nA$ & $\nB$ & $A\cdot B$ & $\overline{A\cdot B}$ & $\nA+\nB$ & $A+B$ & $\overline{A+B}$ & $\nA\cdot\nB$ \\
		\midrule
		0 & 0 & 1 & 1 & 0 & 1 & 1 & 0 & 1 & 1 \\
		0 & 1 & 1 & 0 & 0 & 1 & 1 & 1 & 0 & 0 \\
		1 & 0 & 0 & 1 & 0 & 1 & 1 & 1 & 0 & 0 \\
		1 & 1 & 0 & 0 & 1 & 0 & 0 & 1 & 0 & 0 \\
		\bottomrule
	\end{tabular}
\end{table}

\section{Data Collection}
\label{sec:datacollection10}

\begin{table}[H]
	\centering
	\caption{Experimental Data -- Part A: Inverter (NOT Gate)}\label{tab:partA10}
	\small
	\begin{tabular}{@{}cc@{}}
		\toprule
		Input & Output (LED) \\
		\midrule
		0 & 1 \\
		1 & 0 \\
		\bottomrule
	\end{tabular}
\end{table}

\begin{table}[H]
	\centering
	\caption{Experimental Data -- Part B: AND Gate Configuration}\label{tab:partB10}
	\small
	\begin{tabular}{@{}cccccc@{}}
		\toprule
		Input 1 ($A$) & Input 2 ($B$) & Inverted Input 1 ($\nA$) & Inverted Input 2 ($\nB$) & LED 1 ($A\cdot B$) & LED 2 ($\nA\cdot\nB$) \\
		\midrule
		0 & 0 & 1 & 1 & 0 & 1 \\
		0 & 1 & 1 & 0 & 0 & 0 \\
		1 & 0 & 0 & 1 & 0 & 0 \\
		1 & 1 & 0 & 0 & 1 & 0 \\
		\bottomrule
	\end{tabular}
\end{table}

\begin{table}[H]
	\centering
	\caption{Experimental Data -- Part C: OR Gate Configuration}\label{tab:partC10}
	\small
	\begin{tabular}{@{}cccccc@{}}
		\toprule
		Input 1 ($A$) & Input 2 ($B$) & Inverted Input 1 ($\nA$) & Inverted Input 2 ($\nB$) & LED 1 ($A+B$) & LED 2 ($\nA+\nB$) \\
		\midrule
		0 & 0 & 1 & 1 & 0 & 1 \\
		0 & 1 & 1 & 0 & 1 & 1 \\
		1 & 0 & 0 & 1 & 1 & 1 \\
		1 & 1 & 0 & 0 & 1 & 0 \\
		\bottomrule
	\end{tabular}
\end{table}

\textit{Note: Data transcribed from the instructor-reviewed handwritten laboratory data sheet (Group 04). Logic 1 means the LED is ON and logic 0 means it is OFF.}

\section{Result Analysis}
\label{sec:resultanalysis10}
The inverter results in Table~\ref{tab:partA10} show that the output is always the opposite of the input (0 gives 1, and 1 gives 0), so the circuit behaves as a NOT gate, as expected.

To verify the second theorem, the measured LED~2 column of the AND configuration ($\nA\cdot\nB$) is compared with the complement of the measured $A+B$ column from the OR configuration. For the first theorem, the measured LED~2 column of the OR configuration ($\nA+\nB$) is compared with the complement of the measured $A\cdot B$ column from the AND configuration.

\begin{table}[H]
	\centering
	\caption{Verification of DeMorgan's Theorems from the Measured Data}\label{tab:verification10}
	\small
	\setlength{\tabcolsep}{5pt}
	\begin{tabular}{@{}cc|ccc|ccc@{}}
		\toprule
		$A$ & $B$ & $A+B$ & $\overline{A+B}$ & $\nA\cdot\nB$ & $A\cdot B$ & $\overline{A\cdot B}$ & $\nA+\nB$ \\
		 &  & (measured) & (complement) & (measured) & (measured) & (complement) & (measured) \\
		\midrule
		0 & 0 & 0 & 1 & 1 & 0 & 1 & 1 \\
		0 & 1 & 1 & 0 & 0 & 0 & 1 & 1 \\
		1 & 0 & 1 & 0 & 0 & 0 & 1 & 1 \\
		1 & 1 & 1 & 0 & 0 & 1 & 0 & 0 \\
		\bottomrule
	\end{tabular}
\end{table}

In all four rows the $\overline{A+B}$ column matches the measured $\nA\cdot\nB$ column, and the $\overline{A\cdot B}$ column matches the measured $\nA+\nB$ column. The measured values also agree with the expected truth table (Table~\ref{tab:theory10}) in every row. A few specific observations:
\begin{itemize}[tight]
	\item In the AND configuration, LED~2 is ON only when both inputs are LOW, which is the NOR function, $\overline{A+B}$.
	\item In the OR configuration, LED~2 is OFF only when both inputs are HIGH, which is the NAND function, $\overline{A\cdot B}$.
\end{itemize}
Since the outputs are binary, there is no numerical percentage error to report here: each of the 8 compared outputs either matches or it does not, and all of them matched. Both DeMorgan's Theorems are therefore verified.



\section{Conclusion}
This experiment verified DeMorgan's Theorems using the 7404, 7408 and 7432 logic gate ICs. The inverter produced the complement of its input, the AND configuration showed that $\nA\cdot\nB=\overline{A+B}$, and the OR configuration showed that $\nA+\nB=\overline{A\cdot B}$, with all measured outputs agreeing with the expected truth tables. This confirms that inverting the inputs of an AND gate gives an OR-type function with an inverted output, and vice versa, which is the basis for building any logic function from NAND or NOR gates alone.

\section{References}
\begin{enumerate}[tight]
	\item Mano, M. M., \& Ciletti, M. D. (2013). \textit{Digital Design: With an Introduction to the Verilog HDL} (5th ed.). Pearson. (Chapter 2: Boolean Algebra and Logic Gates).
	\item Floyd, T. L. (2015). \textit{Digital Fundamentals} (11th ed.). Pearson. (Chapter 4: Boolean Algebra and Logic Simplification).
	\item Tocci, R. J., Widmer, N. S., \& Moss, G. L. (2016). \textit{Digital Systems: Principles and Applications} (12th ed.). Pearson. (Chapter 3: Describing Logic Circuits).
\end{enumerate}

\newpage
\section{Appendix: Signed Data Sheet}
\begin{figure}[H]
	\centering
	\includegraphics[height=0.74\textheight]{datasheet/sheet1.png}
	\caption{Instructor-checked laboratory data sheet for this experiment.}
	\label{fig:datasheet10}
\end{figure}
